\section{Discussion}\label{sec:discussion}
\paragraph{What the repaired simulations say, within their limits.}
\begin{itemize}
\item \textbf{Against FP-SGD, two device specifications reach the line on all three tasks} with a chopped Tiki-Taka algorithm:
      the AIHWKIT Li-ECRAM fit and the protonic-inspired scenario, each with c-TTv2 or AGAD (\RVAllThreeFpN{} pairings).
\item \textbf{Against Adam, none does.} On the language model the best analog cell is \RVLmAdamBestGap\,\% worse in bits per
      byte; the analog results inherit FP-SGD's optimiser gap (\RVLmFpSgdVsAdamPct\,\% on this model). Closing it would need
      analog-compatible adaptive updates, which this study does not test.
\item \textbf{Plain SGD at the FP-SGD learning rate} exceeded the margin on every published device and task we ran
      (\RVSgdPubFail{} of \RVSgdPubN{} cells). This is a statement about that protocol, not about SGD after fair tuning.
\item \textbf{The algorithm settings use the device.} On an ideal symmetric 1000-level device the selected TT settings lose
      far more than on the saturating devices, so in this configuration the accumulator's decay comes from the device's
      saturation.
\item \textbf{Width} is not resolved: after per-arm tuning no arm's gap grows detectably beyond the margin over a factor of
      four, endpoint equivalence was established for one CIFAR arm only, the registered Welch test and a paired-by-seed analysis both detect growth for
      another without resolving its size against the margin (growth beyond it is not excluded there or for
      Li-ECRAM TTv2), and the LM's shrinking gap reflects a weaker FP-SGD baseline at the largest width.
\item \textbf{Energy} is a conditional projection set almost entirely by borrowed read efficiencies, and it is reported only
      for cells that pass the stated margin under ideal reads; against Adam on the LM no cell does.
\end{itemize}

\paragraph{Widrow's memistor.} Taken at its inventor's description (thousands of near-linear levels), an optimistic
historical-device scenario attained useful accuracy in this simulator; the pessimistic twin does not. Because its noise,
asymmetry, drift and variability are largely unknown, this does not settle the accuracy of the real 1960 device. Its energy per
level, tens of microjoules, is many orders of magnitude worse than today's devices.

\subsection{Limitations}\label{sec:limits}
\begin{itemize}
\item \textbf{Simulation only.} Every device is a parametric model; the protonic entry is a scenario with estimated noise,
      asymmetry and saturation, and its finite retention is replaced by no drift. A 2026 measurement on BEOL-compatible
      protonic ECRAM~\cite{xu2026} reports 3\,\% cycle-to-cycle and 21\,\% device-to-device variability; we did not refit the
      catalog to it, and we could not check those values against the full text. Estimated parameters are single points, not
      propagated ranges.
\item \textbf{Ideal reads.} Read noise, IR drop, ADC/DAC quantisation and output bounds are not modelled, while the energy
      projection prices finite-precision reads.
\item \textbf{The update surrogate.} Pulses of one batch are aggregated and solved as uniformly interleaved; discrete pulses do
      not commute, per-cell Poisson draws ignore row/column correlations, and the pulse budget clips the TT family's $A$
      updates by up to \RVLostCifarTtLayerMeanMaxPct\,\% in some layers. Against explicit pulse streams
      (Section~\ref{sec:fidelity}) the surrogate meets the registered lines on \RYProbeSurMeetsN{} of \RYProbeN{} update probes
      (point estimates), fails on the most clipped layer and does not reproduce the dependence of cells that share a line; in
      short training comparisons on MNIST and CIFAR's fully connected layers one label loses PASS, two closely spaced cells swap
      order and one failing SGD cell moves by more than 1\,pp. The convolutional layers and the LM's training are not compared,
      and the comparison stream is itself a per-sample approximation of the literal one.
\item \textbf{The reset.} The pulse reset uses a read and at most 62 pulses per weight; it leaves a residual (mean over runs
      \RVResetResidMean{} nominal steps, pooled over all reset events \RVResetResidPooled) and its latency and endurance cost are not modelled.
\item \textbf{Algorithms.} TTv1, TTv2, c-TTv2 and AGAD were implemented from the papers' text; our TTv2 is a local variant
      (multi-pulse transfer, buffer remainder). The settings were selected once, on MNIST, with a mid device.
\item \textbf{Controls.} The ideal-device reference keeps stochastic pulses, bounds and the selected update management, so it
      does not separate the optimiser from the device. A full ladder would run exact digital SGD, a digital TT with an
      explicit accumulator decay, finite pulses and levels on ideal devices, the device imperfections and shared-line pulse
      streams; we ran the first, the ideal 1000-level device, the catalog and, on MNIST and CIFAR's dense layers, the
      per-sample stream.
\item \textbf{Small models and short budgets.} $10^5$ to $10^6$ weights, a 10{,}000-image CIFAR subset, a two-layer byte LM,
      a factor of four in width, fixed training budgets.
\item \textbf{Statistics.} Five seeds (three in the sensitivity analysis and the diagnostic), no multiplicity correction, and best
      cells chosen among many; several cells remain UNRESOLVED.
\item \textbf{Plans.} The plans were committed to our repository before the runs; they are not registered in an external
      registry.
\end{itemize}
